%==============================================================================
% Sjabloon projectvoorstel graduaatsproject
%==============================================================================
% Je gegevens staan in ../projectgegevens.tex. De inhoud van je voorstel
% schrijf je in voorstel-inhoud.tex - dat bestand wordt later ook automatisch
% als bijlage in je rapport opgenomen.

\documentclass{hogent-article}

% Voor een voorstel in het Engels: \documentclass[english]{hogent-article}
% Dat kan enkel na toestemming van je begeleider.

\usepackage{array} 
\addbibresource{voorstel.bib}

%%=============================================================================
%% Projectgegevens - VUL DIT ALS EERSTE IN
%%=============================================================================
%%
%% Dit is de enige plek waar je je naam, je titel, je bedrijf en je begeleiders
%% invult. Het voorstel, het rapport en de poster halen hun gegevens hier op,
%% zodat ze nooit uit elkaar kunnen lopen.
%%
%% Enkel de tekst tussen de accolades aanpassen. Laat de commandonamen staan.

%% ---------- Jij ------------------------------------------------------------
\newcommand{\studentnaam}{Bradley D'Haene}
\newcommand{\studentemail}{bradley.dhaene@student.hogent.be}

%% Wordt gebruikt als bestandsnaam van je PDF's, bijvoorbeeld
%% Familienaam-Voornaam-rapport.pdf. Geen spaties of leestekens gebruiken.
\newcommand{\bestandsnaam}{D'Haene-Bradley}

%% ---------- Je project -----------------------------------------------------
\newcommand{\projecttitel}{Volledige uitwerking RMA website}

%% Een ondertitel is optioneel. Laat leeg als je er geen hebt: {}
\newcommand{\projectondertitel}{}

%% Een of twee zinnen. Wordt gebruikt op de poster en in het voorstel.
\newcommand{\projectpitch}{%
  Het maken van een nieuwe RMA website waar klanten makkelijk formulieren kunnen indienen voor defecte items te retourneren, gebruik makende van C\# & Blazor WASM.
}

%% ---------- Waar je het doet ------------------------------------------------
\newcommand{\bedrijfsnaam}{CMC NV}

%% Je mentor op de werkvloer (het bedrijf)
\newcommand{\mentornaam}{Pieter Eggermont}
\newcommand{\mentoremail}{Pieter.Eggermont@cmcnv.be}

%% Je begeleider van de opleiding (HOGENT)
\newcommand{\begeleidernaam}{Luc Vervoort}
\newcommand{\begeleideremail}{luc.vervoort@hogent.be}

%% ---------- Je repository ---------------------------------------------------
%% De link naar je eigen git-repository. Die hoort zichtbaar in je rapport en
%% op je poster: hij is een beoordeeld onderdeel van je oplevering.
\newcommand{\repolink}{https://github.com/bradleydhaene1994/Graduaatsproject}

%% ---------- Vast voor dit opleidingsonderdeel -------------------------------
%% Deze twee hoef je normaal niet aan te passen.
\newcommand{\opleidingnaam}{Graduaat in het Programmeren}
\newcommand{\academiejaarwaarde}{2026-2027}

%% Zoekwoorden voor het voorstel en de poster (3 tot 5)
\newcommand{\projectkeywords}{Webapp, RMA, Blazor}


%% Opleiding, vak en soort opdracht
\studyprogramme{\opleidingnaam}
\course{Graduaatsproject}
\assignmenttype{Projectvoorstel}

\academicyear{\academiejaarwaarde}

\title{\projecttitel}

\author{\studentnaam}
\email{\studentemail}

%% Je mentor op de werkvloer
\supervisor[Mentor]{\mentornaam{} (\bedrijfsnaam, \href{mailto:\mentoremail}{\mentoremail})}

%% De link naar je repository
\projectrepo{\repolink}

\keywords{\projectkeywords}

\begin{document}

\begin{abstract}
  % Eén samenhangende alinea van vijf tot acht zinnen, die je hele voorstel
  % samenvat: de context bij je bedrijf, het probleem, wat je gaat maken, met
  % welke technologie, en wat het bedrijf eraan heeft. Geen inleiding, geen
  % opsomming.
  Het huidige RMA process bij mijn stagebedrijf is zeer verouderd en vergt zeer veel manuele acties. Ik ga een volledig nieuwe RMA site maken die veel klantvriendelijker is en veel minder manuele 
  handelingen vereist van het personeel. Ik ga hiervoor gebruik maken van C\#, blazor, EF Core SQL Server, Github, Keycloak en Ucloud gehost door het bedrijf uvion. Het uiteindelijke resultaat is een volledig werkend
  geautomatiseerd, klantvriendelijk RMA process via de nieuwe webapp.
\end{abstract}

%%=============================================================================
%% De inhoud van je projectvoorstel
%%=============================================================================
%%
%% Zeven rubrieken, samen hoogstens twee bladzijden. Houd het kort: per
%% rubriek volstaan één tot drie zinnen. Een voorstel is geen rapport - het
%% moet je begeleider in vijf minuten laten beslissen of dit een geschikt
%% graduaatsproject is.
%%
%% Stem je voorstel af met je mentor VOORDAT je het indient. Je mentor
%% bevestigt dat dit is wat het bedrijf bedoelt en dat de tien werkdagen
%% haalbaar zijn.

%---------- 1. Context en probleemstelling ------------------------------------

\section{Context en probleemstelling}%
\label{vst:context}

Ik loop stage bij CMC NV. er is één software team die de software onderhoudt van een aantal producten. Deze producten worden gebruikt in de bedrijfswereld, en als er een probleem is met een product moet er een RMA (return merchandise authorization) aangemaakt worden
om dit te kunnen verwezelijken. Het huidige process is heel verouderd en vergt heel veel manueel werk, also de server waar dit huidige process loopt is heel verouderd, waardoor deze op elk moment offline worden gehaald.
Hier kruipt zeer veel werk in, dus er moet een site worden gemaakt waar alles veel meer geautomiseerd en klantvriendelijk is.

%---------- 2. Doelstelling ---------------------------------------------------

\section{Doelstelling}%
\label{vst:doelstelling}

Op het einde van het project bestaat een website waar men als klant zijnde gemakkelijk een RMA indient. De validatie, bijhouding van de status, etc... zal veel meer geautomiseerd zijn dan het nu is, en er zal veel minder manueel bijgehouden worden.

%---------- 3. Technologieën --------------------------------------------------

\section{Technologieën}%
\label{vst:technologieen}

Er zal gebruik gemaakt worden van de programeertaal C\#, voor de frameworks zal ik gebruik maken van Blazor WASM en EF Core, voor databanken zal ik gebruik maken van Microsoft SQL Server, verder zal ik ook gebruik maken van Github, Keycloak en Uvion Uhost.

Ik veronderstel AI te gebruiken, en dit wordt ook aangemoedigd door mijn stagebedrijf. Ik veronderstel dit vooral in te zetten voor syntax op te zoeken waar ik nog niet vertrouwd mee ben, alsook vooral voor bugfixing.

%---------- 4. Afbakening -----------------------------------------------------

\section{Afbakening}%
\label{vst:afbakening}

Mijn dagelijks stagewerk is meehelpen aan de ontwikkeling van een reporting engine voor rapporten aan te maken via de verschillende applicaties van het bedrijf.

Het project dat ik zelfstanding zou uitwerken is een volledige herwerking van het huidige RMA process, analyse + ontwerp van hoe de website er zou uitzien + volledige nieuwe codebase uitschrijven voor nieuwe website.

%---------- 5. Verwachte meerwaarde -------------------------------------------

\section{Verwachte meerwaarde}%
\label{vst:meerwaarde}

Het RMA process zal veel vlotter en veel geautomatiseerder verlopen. Doordat da status van alle RMA's niet meer manueel moet bijgehouden worden, zal dit voor veel tijdswinst zorgen en tot minder fouten lijden.

%---------- 6. Planning en werkbelasting --------------------------------------

\section{Planning en werkbelasting}%
\label{vst:planning}

% Het graduaatsproject is ongeveer tien werkdagen, samen ongeveer 100 uren.
% Verdeel die over de fasen van je project en geef per fase aan wat je oplevert.
% Vermeld ook of je die dagen gefocust opneemt of gespreid over de stageperiode,
% en dat je dit met je mentor afgesproken hebt.

\begin{table}[htbp]
  \centering
  \begin{tabular}{llr}
    \toprule
    \textbf{Fase} & \textbf{Wat lever je op} & \textbf{Uren} \\
    \midrule
    Analyse en vereisten & overzicht van de vereisten   & 5  \\
    Technologiekeuze     & verantwoorde keuze           & 5  \\
    Ontwerp              & architectuur en gegevensmodel& 20  \\
    Realisatie           & werkende oplossing           & 50  \\
    Testen en validatie  & testverslag                  & 10  \\
    Rapport, poster en demo & de op te leveren documenten  & 10  \\
    \midrule
    \textbf{Totaal}           &                              & \textbf{100} \\
    \bottomrule
  \end{tabular}
  \caption[Planning en werkbelasting]{\label{tab:vst-planning}%
    }
\end{table}

%---------- 7. Risico's -------------------------------------------------------

\section{Risico's}%
\label{vst:risicos}

1. Blazor is een nieuwe technologie waar ik nog niet vertrouwd mee ben: Meer tijd plannen om documentatie te doornemen. 
2. Problemen met hosten via Uvion. Een bepaalde zekerheid inbouwen in de planning waardoor problemen met hosting tijdig kunnen worden ontdekt en opgelost.


\printbibliography[heading=bibintoc]

\end{document}
